% =========================================================================
% secciones_zh/sec01_introduction.tex — 第I节，中文版（精修版）
% =========================================================================
\section{引言}
\label{sec:intro}

孤立互联电网处于电力工程设计空间的极端位置。岛屿无法从邻近电网输入惯量、无功支撑或紧急输送功率：每一兆瓦的备用容量、每一伏安的无功补偿和每一兆赫的阻尼能力，都必须在消费它的同一同步覆盖区内生产、调度和防御。多米尼加共和国是一个典型案例：一个约有一千一百万人口的单一同步岛，无海底或陆上互联，全国百分之四十以上的负荷集中在奥萨马（Ozama）都市区，惯量锚点由蓬塔卡塔利纳（Punta Catalina）燃煤电厂主导，且可再生能源建设的瞬时光伏渗透率在运行遥测中已达到需求的36.6\%\cite{ocseni}。失去发电--负荷平衡所造成的全国性后果并非假设：2025年11月与2026年2月记录在案的全国性停电事件\cite{ocseni}曾以整个领土的规模动员应急响应，并重新引发了关于频率稳定性、黑启动能力与输电资产可观测性的公共政策辩论\cite{blackout2004}。这种脆弱性的物理机理如今已是经典\cite{kundur1994}——系统惯量下降会劣化频率变化率并收窄变流器主导电网的可控包络\cite{tielens2016}、\cite{milano2018}、\cite{dorfler2013}、\cite{hatziargyriou2021}——但可用于研究它的实证工具仍然出奇地原始。

本领域的基准测试基础设施大多是合成的。至今仍主导博士训练与对比研究的经典测试系统——WSCC 9节点系统以及IEEE 14、30、57、118节点系统——都源自二十世纪六十年代初的制图台式抽象：电抗是圆整的学术数字，负荷是固定的功率三角形，没有地理，没有气候，数据集中没有任何元件曾跳闸、检修、付费或遭受过飓风。另一个极端上，电力公司级的电网案例保存在商业仿真平台的专有格式中，受保密约束，很少以允许独立重新执行的形态发布。开源运动提供了宝贵的基础设施——MATPOWER算例库\cite{matpower2011}与AC-OPF基准算例的电力网格库（Power Grid Library）\cite{pglib2019}等——但其算例几乎无一例外是\emph{合成或匿名的}：它们是为"难"而生的电网，而不是真实存在的电网。在玩具与机密之间几乎一无所有：一个开放的、完整的、实测的真实国家输电系统模型，其溯源性足以支持独立验证——这是加勒比海岛屿世界（乃至整个岛屿世界）当前所缺乏的公共品。

本文以\textbf{SENI-142}弥合这一鸿沟：一个围绕三项承诺构建的开放基准。第一，\textbf{对实测现实的忠实}：网络模型不是合成的——它源自国家输电资产，并以87,360个连续运行小时（2016--2025年）的系统运营商遥测数据逐机组、逐小时地进行验证，使每个仿真状态都对应一个存档的运行状态。第二，\textbf{对尺度的结构诚实}：求解的是全国完整系统——全部输电母线、全部干线、全部变压器——而非其缩减片区，且求解方法以定理而非愿望的形式给出其精确性与复杂度：网络分块（diakoptics）传统\cite{kron1963}、\cite{happ1971}下的分层边界-坐标列式（HBCF），经舒尔补\cite{zhang2005}凝聚，并在国家领土规划法（第368-22号法律）\cite{ley368}规定的十个领土规划区上分解因式。第三，\textbf{确定性可复现}：每个工件都通过具有失效即闭（fail-closed）校验门级联、SHA-512分离式封存与原子化晋升的四阶段数字孪生生命周期，使本文报告的任何数字都可以从逐位相同、经密码学证明的软件状态重新生成。

贡献有五项，且每项均可证伪：
\begin{itemize}
\item \textbf{C1 —— 基准。}多米尼加共和国完整输电网络的开放、地理配准、符合标准的模型：约112座物理变电站之上的142个节点，135条输电线路，58台电力变压器，345、230与138~kV三个电压层，69~kV subtransmission层经保持注入的等效化简；随附拓扑、参数、溯源性元数据与验证工具，通过IEEE DataPort与GitHub开放发布。
\item \textbf{C2 —— 列式。}将完整岛屿电网网络分块划分为十个主权领土微电网——即第368-22号法律的规划区——并聚合为四个电气宏域、边界集最小的HBCF划分，以及其舒尔补凝聚的精确性证明：复杂度 $O(\sum_k N_k^3 + N_B^3)$，对比整体稠密分解的 $O(N^3)$。
\item \textbf{C3 —— 确定性数字孪生。}四阶段生命周期（带十四道失效即闭校验门级联的POSIX命名空间密封飞地；针对215个预想事故目录的动态科研测试台 $N-1$ 压力测试；SHA-512分离式封存的黄金主版本；经 \texttt{renameat2} 的原子化内核晋升，回归量为零），将以链式凭证把仿真遥测转化为证据。
\item \textbf{C4 —— 多向量灵活性。}将电转气氢能（电解、700~bar存储、燃料电池再注入）、764~MW可逆抽水蓄能与基于ADM1的生物甲烷化耦合到多时段AC-OPF，给出完整的Karush--Kuhn--Tucker（KKT）条件体系，并以分区CAPEX/OPEX平衡证明加以验证。
\item \textbf{C5 —— 实测硅上性能。}在消费级硬件上以2.65~ms、473~MiB常驻内存与 $\delta = 6.35 \times 10^{-12}$~pu 的相量边界残差求解每个全国小时级交流状态——快于60~Hz下8.33~ms的半周期实时阈值——并采用24小时块自助法验证协议。
\end{itemize}

本文余下部分的组织如下。第~\ref{sec:methodology}节阐述方法论：不变列式（图模型、69~kV化简、HBCF划分、舒尔补凝聚、数据协议、确定性生命周期、无功限值正则化、多时段AC-OPF）与在多米尼加案例上按第368-22号法律实例化的领土条款。第~\ref{sec:topology}节记录基准拓扑：从112座变电站的物理资产到142节点模型与十个领土微电网。第~\ref{sec:multivector}节展开多向量灵活性星座及其分区经济性验证。第~\ref{sec:benchmarking}节报告硅上性能、精度与预想事故韧性。第~\ref{sec:discussion}节讨论政策启示与局限，第~\ref{sec:conclusion}节以开放发布计划与四条研究路线作结。
